\part{Aktuelle Aktivitäten}

\begin{frame}{Aktuelle Aktivitäten}
  Hauptaktivität ist aktuell die \alert{Aufnahme} der \alert{\texttt{MQTTBroker/MQTTIntegrator}} Referenzapplikationen in \alert{\texttt{OpenWRT}}, eine auf \alert{Linux} basierende \alert{Open-Source Firmware} für (Wi-Fi) Router.

  \color{fh-ooe-red}{Nahezu jede proprietäre Router-Firmware für Heim-Router basiert auf \texttt{OpenWRT}.}
  \begin{itemize}
    \item Zwei Libraries waren noch nicht Bestandteil von \texttt{OpenWRT}
    \begin{itemize}
      \item \alert{\texttt{Easylogging++}} (Logging Library)
      \item \alert{\texttt{libucontext}} (kooperatives Multithreading)
      \begin{itemize}
        \item ARM: Bug Fix
        \item Aarch64: Bug Fix
      \end{itemize}
      OpenWRT verwendet die \alert{\texttt{musl}} statt der \alert{\texttt{GNU-C}} Library. \texttt{muls} stellt im Gegensatz zur \texttt{GNU-C} Library die Systemaufrufe zum kooperativen Multithreading nicht zur verfügung, da diese seit \texttt{POSIX.1-2008} aus dem Standard genommen wurden.
    \end{itemize}
    \item Update auf \alert{\texttt{upstream}} bei zwei Packeten von \texttt{OpenWRT}
    \begin{itemize}
      \item \alert{\texttt{File}}
      \item \alert{\texttt{nlohmann\_json}}
    \end{itemize}
    \item Fix \alert{\texttt{segfault}} im \alert{asynchronen API} von \alert{\texttt{libmariadb}} durch hinzufügen des linkens gegen \alert{\texttt{libucontext}} (noch nicht beendet)
    \item Und letztendlich muss \alert{\texttt{SNode.C}} selbst in OpenWRT aufgenommen werden, da die MQTT-Applicationen auf SNode.C aufsetzen.
  \end{itemize}
\end{frame}
